\documentclass[11pt]{article}
\usepackage[a4paper,margin=1.9cm]{geometry}
\usepackage[T1]{fontenc}
\usepackage{textcomp}
\usepackage[spanish]{babel}
\usepackage{amsmath,amssymb}
\usepackage{lmodern}
\renewcommand{\familydefault}{\sfdefault}
\usepackage{enumitem}

\newcommand{\vect}[2]{\begin{pmatrix}#1\\#2\end{pmatrix}}
\newcommand{\vectiii}[3]{\begin{pmatrix}#1\\#2\\#3\end{pmatrix}}
\usepackage{tikz}
\usetikzlibrary{calc,arrows.meta,patterns,decorations.pathmorphing}
\usepackage[table]{xcolor}
\usepackage{multicol}
\usepackage{parskip}

\definecolor{unalblue}{RGB}{31,73,124}

\newlist{opciones}{enumerate}{1}
\setlist[opciones]{label=(\alph*),itemsep=6pt,topsep=6pt,leftmargin=1.8em,font=\normalfont}
\setlist[enumerate,1]{itemsep=1.4em,topsep=1em}

\newcommand{\examfooter}{%
  \vfill
  \rule{\linewidth}{0.4pt}\\[1pt]
  {\scriptsize\textit{Transcripción digital --- MathUNAL}\hfill\textcolor{unalblue}{\textbf{\thepage}}}
}
\pagestyle{empty}

\newcommand{\nota}[1]{{\small\itshape\color{unalblue!70!black} #1}}

\begin{document}
\noindent
\begin{minipage}[t]{0.6\linewidth}
{\small\bfseries Geometría Vectorial y Analítica --- Parcial 3}\\
{\small Semestre 2016--02}
\end{minipage}%
\hfill
\begin{minipage}[t]{0.35\linewidth}
\raggedleft\small Escuela de Matemáticas. Universidad Nacional de Colombia, Sede Medellín.
\end{minipage}\\[2pt]
\rule{\linewidth}{0.6pt}

\begin{center}
\textbf{Tercer Examen Parcial de Geometría}\\
21 de noviembre del 2016
\end{center}

\vspace{6pt}
\textbf{Puntaje. Sólo para uso Oficial}

\begin{center}
\begin{tabular}{|c|c|c|c|c|}
\hline
1 (/30) & 2 (/30) & 3-4 (/40) & Total & Nota \\
\hline
 & & & & \\
\hline
\end{tabular}
\end{center}

\vspace{10pt}
\textbf{Instrucciones:} La duración del examen es de 1 hora y 50 minutos. El examen consta de cuatro preguntas en tres hojas impresas por ambos lados, antes de empezar verifique que su examen esté \textbf{completo} y que sus \textbf{datos} sean \textbf{correctos}. No \textbf{desengrape} su examen ni \textbf{añada} otras grapas o su examen será \textbf{anulado}. En las preguntas con procedimiento \textbf{justifique} sus respuestas en los \textbf{espacios asignados}. No está permitido hacer preguntas, el uso de calculadoras, sacar hojas en blanco ni ningún tipo de apuntes durante el examen. Por favor verifique que su celular esté \textbf{apagado}.

\vspace{6pt}
\nota{En cada uno de los literales de los puntos 1 y 2, debe presentar detalladamente el procedimiento para llegar a la solución. NOTA: Respuestas sin procedimiento no se tendrán en cuenta.}

\begin{enumerate}
\item{}[30pt] En la figura se observa un cubo, \textbf{cuya arista mide 1}. El plano $\mathcal{P}$ contiene las tres esquinas $P,Q,R$ del cubo.

\begin{center}
\begin{tikzpicture}[scale=1.6]
\coordinate (O) at (0.9,0.55);
\coordinate (R) at (0,0.15);
\coordinate (P) at (1.9,0.15);
\coordinate (S) at (0.95,0);
\coordinate (Otop) at (0.9,1.75);
\coordinate (Rtop) at (0,1.35);
\coordinate (Ptop) at (1.9,1.35);
\coordinate (Q) at (0.95,1.2);
\draw[dotted] (O) -- (R);
\draw[dotted] (O) -- (P);
\draw[dotted] (O) -- (Otop);
\draw[thick] (R) -- (S) -- (P);
\draw[thick] (R) -- (Rtop) -- (Otop) -- (Ptop) -- (P);
\draw[thick] (Rtop) -- (Q) -- (Ptop);
\draw[thick] (S) -- (Q);
\draw[dashed,thick] (Q) -- (R);
\draw[dashed,thick] (Q) -- (O);
\draw[dashed,thick] (Q) -- (P);
\draw[-{Latex}] (Otop) -- ($(Otop)+(0,0.35)$) node[above]{$z$};
\draw[-{Latex}] (O) -- ($(R)+(-0.15,-0.15)$) node[left]{$x$};
\draw[-{Latex}] (O) -- ($(P)+(0.15,-0.05)$) node[right]{$y$};
\fill (Q) circle (0.7pt); \node[above] at (Q) {$Q$};
\fill (R) circle (0.7pt); \node[left] at (R) {$R$};
\fill (P) circle (0.7pt); \node[right] at (P) {$P$};
\fill (S) circle (0.7pt); \node[below] at (S) {$S$};
\node[above left,font=\small] at (O) {$O$};
\node[right,font=\small] at (1.1,1) {$\mathcal{P}$};
\end{tikzpicture}
\end{center}

\begin{enumerate}[label=(\alph*)]
\item{}[5pt] Encuentre las coordenadas de los puntos $P$, $Q$, $R$ y $S$ respectivamente.
\item{}[9pt] Encuentre una ecuación general del plano $\mathcal{P}$.
\item{}[8pt] Encuentre la distancia del punto $S$ al plano $\mathcal{P}$.
\item{}[8pt] Encuentre el área del triángulo $\triangle PQR$.
\end{enumerate}

\item{}[30pt] Considere el conjunto de puntos $P=\vect{x}{y}\in\mathbb{R}^2$ para los cuales la distancia de $P$ al punto $B=\vect{-5}{0}$ es igual al doble de la distancia de $P$ a la recta $\mathcal{L}: x+2=0$.

\begin{enumerate}[label=(\alph*)]
\item{}[10pt] Describa mediante una ecuación en las variables $x,y$, el conjunto de puntos $P=\vect{x}{y}$ del plano que satisfacen la condición descrita en el enunciado.
\item{}[12pt] Compruebe que los puntos $P$ forman una hipérbola y escriba su ecuación en forma estándar.
\item{}[8pt] Encuentre el centro de la hipérbola y la distancia entre sus vértices.
\end{enumerate}
\end{enumerate}

\vspace{6pt}
\nota{En las preguntas 3 y 4 complete los espacios en blanco o elija la (las) opción (opciones) correcta (correctas) según el caso. Tenga presente que los datos en cada literal son independientes y pueden variar. NOTA: En esta sección se califica únicamente la respuesta y no se tiene en cuenta el procedimiento.}

\begin{enumerate}
\setcounter{enumi}{2}
\item{}[20pt] En la gráfica se observan tres parábolas $\mathcal{P}_1,\mathcal{P}_2,\mathcal{P}_3$ con vértice en el origen. Las tres rectas horizontales $\mathcal{L},\mathcal{L}',\mathcal{L}''$ son directrices de las parábolas, no necesariamente en orden. Responda a las siguientes preguntas.

\begin{center}
\begin{tikzpicture}[scale=0.8]
\draw (-4,-2.6) -- (-4,2.6); \node[below] at (-4,-2.6) {$\mathcal{L}'$};
\draw (-3,-2.6) -- (-3,2.6); \node[below] at (-3,-2.6) {$\mathcal{L}$};
\draw (-2.2,-2.6) -- (-2.2,2.6); \node[below] at (-2.2,-2.6) {$\mathcal{L}''$};
\draw[-{Latex}] (-4.3,0) -- (2.6,0) node[right]{$x$};
\draw[-{Latex}] (0,-2.7) -- (0,2.7) node[above]{};
\draw[thick,domain=-2.3:2.3,smooth,variable=\t] plot ({0.45*\t*\t},{\t}) node[right]{$\mathcal{P}_2$};
\draw[thick,domain=-2.1:2.1,smooth,variable=\t] plot ({0.65*\t*\t},{\t}) node[right]{$\mathcal{P}_1$};
\draw[thick,domain=-1.75:1.75,smooth,variable=\t] plot ({0.9*\t*\t},{\t}) node[right]{$\mathcal{P}_3$};
\fill (-4,1) circle (1.3pt); \node[left] at (-4,1) {$Y$};
\fill (-2,1) circle (1.3pt); \node[above] at (-2,1) {$X$};
\draw[dashed] (-4,1) -- (0.9,1);
\fill (1.4,0) circle (1.3pt); \node[below] at (1.4,0) {$F_1$};
\end{tikzpicture}
\end{center}

\begin{enumerate}[label=(\alph*)]
\item{}[5pt] Si el foco de la parábola $\mathcal{P}_3$ tiene coordenadas $F_3=\vect{5}{0}$ entonces su directriz tiene la ecuación: Directriz: \underline{\hspace{4cm}}.
\item{}[5pt] Si el lado recto de la parábola $\mathcal{P}_2$ mide 8 unidades, entonces su foco tiene coordenadas $F_2=$ \underline{\hspace{3cm}}.
\item{}[5pt] Suponga que el punto $X=\vect{1/2}{4}$ está en la parábola $\mathcal{P}_1$, cuya directriz es $\mathcal{L}'$. Además el punto $Y\in\mathcal{L}'$ satisface $\operatorname{dist}(X,Y)=\operatorname{dist}(X,\mathcal{L}')$. Si el foco de $\mathcal{P}_1$ es $F_1=\vect{7/2}{0}$, entonces $Y$ tiene coordenadas $Y=$ \underline{\hspace{3cm}}.
\item{}[5pt] ¿Cuál de las tres rectas horizontales $\mathcal{L},\mathcal{L}'$ o $\mathcal{L}''$ es la directriz de $\mathcal{P}_2$? Respuesta: \underline{\hspace{2.5cm}}.
\end{enumerate}

\item{}[20pt] En la figura se representan los puntos $P,Q,R,S$ y $O=\vectiii{0}{0}{0}$ de $\mathbb{R}^3$. También se observan planos horizontales en las alturas indicadas, los cuadrados chicos miden $1\times 1$.

\begin{center}
\begin{tikzpicture}[scale=0.75]
\foreach \h in {0,1,2,3,4}{
  \draw[gray!50,thin] (-1.5,-\h*0.4) -- (2,-\h*0.4+1) -- (5,-\h*0.4) -- (1.5,-\h*0.4-1) -- cycle;
}
\draw[-{Latex}] (1.75,-2.1) -- (1.75,3.4) node[above]{$z$};
\draw[-{Latex}] (1.75,-2) -- (-1.9,-2.5) node[left]{$y$};
\draw[-{Latex}] (1.75,-2) -- (5.3,-2.65) node[right]{$x$};
\coordinate (O) at (1.75,-2);
\coordinate (S) at (0.3,-2.3);
\coordinate (P) at (3.2,-2.4);
\coordinate (Q) at (1.4,-0.6);
\coordinate (R) at (2.1,1.1);
\draw[-{Latex},thick] (O) -- (S) node[below]{$S$};
\draw[-{Latex},thick] (O) -- (P) node[below right]{$P$};
\draw[-{Latex},thick] (O) -- (Q) node[right]{$Q$};
\draw[-{Latex},thick] (O) -- (R) node[above right]{$R$};
\node[below left,font=\small] at (O) {$O$};
\end{tikzpicture}
\end{center}

\begin{enumerate}[label=(\alph*)]
\item{}[5pt] Si $\mathcal{P}=\{tQ+sR: t,s\in\mathbb{R}\}$, entonces un vector normal $N$ del plano $\mathcal{P}$ es $N=$ \underline{\hspace{3.5cm}}.
\item{}[5pt] Una ecuación vectorial de la recta $\mathcal{L}$ que pasa por los puntos $R$ y $S$ es $\mathcal{L}=$ \underline{\hspace{3.5cm}}.
\item{}[5pt] Si $\mathcal{L}_1=\{tS+3E_3: t\in\mathbb{R}\}$ y $\mathcal{L}_2=\{tP: t\in\mathbb{R}\}$ entonces, la distancia entre estas rectas es $\operatorname{dist}(\mathcal{L}_1,\mathcal{L}_2)=$ \underline{\hspace{3.5cm}}.
\item{}[5pt] El volumen del paralelepípedo determinado por $P$, $R$ y $S$ es Volumen $=$ \underline{\hspace{3.5cm}}.
\end{enumerate}

\end{enumerate}

\examfooter
\end{document}
